% sections/1_intro.tex
\section{Introducción}

El modelo de Ising bidimensional constituye uno de los paradigmas fundamentales en la física estadística para el estudio de transiciones de fase y fenómenos críticos. Este sistema describe la dinámica de espines magnéticos interactuantes en una red discreta, permitiendo modelar el ferromagnetismo a nivel microscópico. 


El objetivo principal de este trabajo es simular computacionalmente el modelo de Ising 2D en ausencia de campo magnético externo, enfocando el análisis en la transición de fase continua que ocurre a la temperatura crítica ($T_c$). El alcance computacional abarca el estudio de la energía, la capacidad calorífica, la magnetización (como parámetro de orden del sistema) y la susceptibilidad magnética. Adicionalmente, se busca validar mediante un análisis de escalado de tamaño finito (\textit{finite-size scaling}) las divergencias teóricas de la capacidad calorífica y la susceptibilidad en el límite termodinámico. 
Posteriormente, el estudio se amplía mediante la introducción de un campo magnético externo, lo que permite explorar el régimen de rotura explícita de simetría y caracterizar la respuesta del sistema a través de su ciclo de histéresis magnética. Por último, se introduce un enfoque basado en aprendizaje automático (\textit{Machine Learning}), empleando un modelo de regresión logística para automatizar la clasificación del sistema en sus fases ordenada (ferromagnética) y desordenada (paramagnética) a partir de sus observables termodinámicos.